% concatenated from shiu_for_journal.tex
\documentclass{amsart}
\usepackage{enumitem}
\usepackage{latexsym}
\usepackage{amsmath,amsfonts,amsthm, amssymb}
\input{xy}
\xyoption{all}


\newtheorem{theorem}{Theorem}[section]
\newtheorem{lemma}[theorem]{Lemma}
\newtheorem{corollary}[theorem]{Corollary}
\newtheorem{varthm1}{Main Theorem}
\newtheorem{varthm2}{Main Theorem 2}
\newtheorem*{varthmdfold}{Theorem {\ref{dfold}}}
\newtheorem*{vtsmooththm}{Theorem {\ref{smooththm}}}
\newtheorem*{varthmHilcor}{Corollary {\ref{Hilcor}}}
\newtheorem{conjtd}{Conjecture}
\newtheorem{conjth}{Conjecture/Theorem}
\newtheorem{ques}{Question}
\newtheorem{cond}{Condition}
\newtheorem*{condi}{Condition i}
\newtheorem*{condii}{Condition ii}
\begin{document}


\title{A Shiu Theorem for Larger and Smoother Functions}
\author[T. Wright]{Thomas Wright}
\begin{abstract}
In this paper, we broaden Shiu's Brun-Titchmarsh theorem to allow for functions that are larger and/or smooth-supported.  In particular, let $f$ be a nonnegative multiplicative function.  We prove that if there exists a $\beta<1$ such that $f(p^l)\ll (\log\log x)^{l\beta}$ for every prime $p$ and every $l>1$, and if $f(n)\ll \max\{n^\epsilon,(\log x)^\epsilon\}$ for every $\epsilon>0$, then
$$\sum_{\substack{x\leq n\leq x+y \\ n\equiv a\pmod k}}f(n)\ll \frac{y}{\phi(k)(\log x)^{1-\epsilon_0}}\exp\left(\sum_{\substack{p\leq x \\ p\nmid k}}\frac{f(p)}{p}\right)$$
for every $\epsilon_0>0$, where $x$, $y$, and $k$ are as they were in Shiu's original paper and $(a,k)=1$.  Moreover, we prove that if $f$ is a $Q$-smooth-supported function then there exists a constant $C$ for which
$$\sum_{\substack{x\leq n\leq x+y \\ n\equiv a\pmod k}}f(n)\ll \frac{y}{\phi(k)(\log x)^{1-\epsilon_0}}\exp\left(\sum_{\substack{p\leq x \\ p\nmid k}}\frac{f(p)}{p}\right)\rho(u)^C,$$
where $u=\frac{\log x}{\log Q}$, $\rho$ is the Dickman-de Bruijn function, and $C$ depends on whether we choose the bound of $f(p^l)\leq A_1^l$ or $f(p^l)\ll (\log\log x)^{l\beta}$.

We also give applications to both the divisor function to large powers and to smooth numbers in short intervals.
\end{abstract}
\maketitle


\section{Introduction}

One of the most widely-used bounds in the theory of multiplicative functions was proven by Peter Shiu in 1980 \cite{Sh}, who proved a Brun-Titchmarsh-type result for these functions.  We begin by stating this result below.

Let $\mathcal M$ denote the set of nonnegative multiplicative functions $f$ such that the following two conditions hold.  

\begin{condi}
There exists a constant $A_1$ such that for any prime $p$ and any natural number $l$,
$$f(p^l)\leq A_1^l.$$
\end{condi}
\begin{condii}
For any $\epsilon>0$ there exists a constant $A_2=A_2(\epsilon)$ such that for any $n$,
$$f(n)\leq A_2n^\epsilon.$$
\end{condii}


Shiu's theorem is then the following.

\begin{theorem}[Shiu, 1980]\label{shiu}
Let $f\in \mathcal M$.  Let $0<\alpha<\frac 12$ and $0<\kappa<\frac 12$, and let $x^\kappa\leq y\leq x$ and $k<y^{1-\alpha}$.  Moreover, let $a$ be such that $(a,k)=1$.  Then for all sufficiently large $x$,
$$\sum_{\substack{x\leq n<x+y \\ n\equiv a\pmod k}}f(n)\ll \frac{y}{\phi(k)\log x}\exp\left(\sum_{\substack{p\leq x\\ p\nmid k}}\frac{f(p)}{p}\right).$$
Here, the term "sufficiently large" depends only on the choices of $\alpha$, $\kappa$, and $\epsilon$ and the constants $A_1$ and $A_2(\epsilon)$.
\end{theorem}


However, there are a few limitations to this theorem that we wish to address in the current paper.  First, this theorem cannot say anything when $f(p)$ is unbounded.  Second, the result is limited in the case where $f$ is a smooth-supported function, since smooth numbers are generally far more sparse than this bound would indicate.  In this paper, we expand this theorem to handle both such cases.

\section{Introduction: Larger Functions}

Let us begin with the idea of letting $f(p)$ go to infinity.  We modify Shiu's condition first condition in the following way:  
\begin{cond}\label{Condition1}
For some fixed $\beta>0$, there exists an $A_1$ such that for any prime $p$,
\begin{gather}\label{Cond1}
f(p^l)\leq A_1^l (\log\log x)^{\beta l}.
\end{gather}
\end{cond}
For the second condition, we would like to use something similar to Condition ii above.  However, this is too restrictive for our purposes, since Condition 1 would allow $f(2)$ to be as big as $(\log\log x)^{\beta}$, which would contradict the requirement that $f(2)\ll 2^\epsilon$. 

To counter this, we add an extra term to handle the case of $p$ small:
\begin{cond}\label{Condition2}
For any $\epsilon>0$, there exist $A_2=A_2(\epsilon)$ and $A_3=A_3(\epsilon)$ such that for every $n$,
\begin{gather}\label{Cond2}
f(n)\leq \max\{A_2 x^\epsilon,A_3(\log x)^\epsilon\}.
\end{gather}
\end{cond}
Since Condition 1 depends on $\beta$, we will denote by $\mathcal M(\beta)$ the set of nonnegative multiplicative functions that satisfy Condition 1 for a given $\beta$ and Condition 2 as above.
\begin{theorem}\label{Maintheorem}
Let $x$ be large.  Let $0<\alpha<\frac 12$ and $0<\kappa<\frac 12$, and let $x^\kappa\leq y\leq x$ and $k<y^{1-\alpha}$.  Moreover, let $f\in \mathcal M(\beta)$ for some fixed $\beta<\frac{\alpha\kappa}{41}$.  Then for any $\epsilon_0>0$,
$$\sum_{\substack{x\leq n\leq x+y \\ n\equiv a\pmod k}}f(n)\ll \frac{y}{\phi(k)(\log x)^{1-\epsilon_0}}\exp\left(\sum_{\substack{p\leq x \\ p\nmid k}}\frac{f(p)}{p}\right).$$
\end{theorem}

We remark that the $\frac{1}{(\log x)^{1-\epsilon_0}}$ in the theorem above could be improved to $\frac{e^{(\log\log x)^{1-\epsilon_0}}}{\log x}$ if we tightened Condition 2 to 
\begin{gather*}
f(n)\leq \max\{A_2 x^\epsilon,A_3e^{(\log\log x)^{1-\gamma}}\}
\end{gather*}
for some $0<\gamma<1$.  


Further tightening of Condition 2 to smaller orders of magnitude may well lead to further improvement.  However, it appears that this method could not fully recover the $1/\log x$ that appears in Shiu's original argument as long as $f(p)$ remains an unbounded function.

\section{Introduction: Smooth-Supported Functions}


Next, we examine the case where $f$ is smooth-supported function.  For a given $Q\leq x$, define
$$u=\frac{\log x}{\log Q},$$
let $\psi(x,Q)$ denote the number of $Q$-smooth numbers up to $x$, and let $\mathcal S(Q)$ denote the $Q$-smooth numbers.  It is known that when $Q\geq \exp\left((\log\log x)^{\frac 53+\epsilon}\right)$, a value for $\psi(x,Q)$ is given by the formula
\begin{gather}\label{Di}\psi(x,Q)=x\rho(u)(1+o(1)),\end{gather}
where $\rho$ is the Dickman-de Bruijin function.  It is believed that this holds for all $Q\geq (\log x)^{2+\epsilon}$, and indeed this is known to be true under the Riemann Hypothesis, and it is known that 
\begin{gather}\label{Di}\psi(x,Q)\leq x\rho(u)^{1+o(1)},\end{gather}
uniformly for $x>y\geq 2$ \cite{Hi}.  Importantly, 
$$\rho(u)=u^{-u+o(1)}$$
holds for all $u$.  

On a short interval, we have a 1985 result of Hildebrandt \cite{Hi85}, which states that if $Q$ is sufficiently large and $x>y>Q$ then
\begin{gather}\label{Hild}\psi(x+y,Q)-\psi(x,Q)\leq \psi(y,Q).
\end{gather}


On the other hand, Shiu's theorem with $k=1$ gives a bound of 
$$\sum_{\substack{x\leq n\leq x+y \\ n\in \mathcal S(Q)}}1\ll \frac{y\log Q}{\log x}.$$
Shiu's bound is clearly not sharp here, as $\frac{y\log Q}{\log x}$ may be much larger than $\psi(y,Q)$.  


%Hence if $Q$ is less than, say, $\leq e^{\frac{\log x\log\log\log x}{2\log\log x}}$ then 
%$$\rho(u)\ll \frac{1+o(1)}{\log^2 x}$$
Here, we bring the Shiu bound a bit closer to expectations for any multiplicative function $f$.  For ease of notation, we will define $\mathcal M_Q$ to be the set of $Q$-smooth supported nonnegative functions $f$ in $\mathcal M$, and we will define $\mathcal M_Q(\beta)$ analogously.

\begin{theorem}\label{MnThm2}
Let $x$, $y$, $\alpha$, $\kappa$, and $k$ be as in Theorem \ref{Maintheorem}.  There exists a constant $C_1$ such that if $f\in \mathcal M_Q$ then
$$\sum_{\substack{x\leq n<x+y \\ n\equiv a\pmod k}}f(n)\ll \rho(u)^{C_1}\frac{y}{\phi(k)\log x}\left(\sum_{\substack{p\leq Q \\ p\nmid k}}\frac{f(p)}{p}\right).$$
Moreover, let $Q\leq e^{\frac{\log x\log\log\log x}{\log\log x}}.$  Then there exists a $\beta<\frac{\alpha\kappa}{41}$ and a constant $C_2$ such that if $f\in \mathcal M_Q(\beta)$ then for any $\epsilon>0$,
$$\sum_{\substack{x\leq n<x+y \\ n\equiv a\pmod k}}f(n)\ll \rho(u)^{C_2}\frac{y}{\phi(k)(\log x)^{1-\epsilon}}\left(\sum_{\substack{p\leq Q \\ p\nmid k}}\frac{f(p)}{p}\right).$$
\end{theorem}
It will suffice to take 
$$C_1=\frac{\alpha\kappa}{41}$$ 
and
$$C_2=\frac{5}{656}\left(\frac{\alpha\kappa \left(1-\frac{\alpha \kappa}{20}\right)}{20}\right).$$

We note here that while we do not put a lower bound on $Q$, Theorem \ref{MnThm2} ceases to be non-trivial when $Q=(\log y)^K$ for a fixed $K$.  This is because such a $Q$ would have $\rho\left(\frac{\log y}{\log Q}\right)\ll y^{1-\frac{1}{K}+\epsilon}$, and hence using (\ref{Hild}) with the trivial bound of $f(n)\ll y^\epsilon$ gives a bound of
$$\sum_{\substack{x\leq n<x+y \\ n\equiv a\pmod k}}f(n)\ll y^\epsilon\sum_{\substack{x\leq n<x+y \\ n\equiv a\pmod k}}1\ll y^{1-\frac 1K+\epsilon},$$
while our bound would give
$$\sum_{\substack{x\leq n<x+y \\ n\equiv a\pmod k}}f(n)\ll \frac{y^{1-\frac{C_i}K}}{\phi(k)(\log x)^{1-\epsilon}}.$$
%In the case where $f(n)$ is the characteristic function on $Q$-smooth numbers, we can sharpen the bound for the %constant.
%\begin{theorem}
%Let $y=x^\eta$ for some fixed $\eta>0$.  Then for any $\nu>0$,
%$$\sum_{\substack{x\leq n\leq x+y \\ n\equiv a\pmod k \\ n\in \mathcal S(Q)}}1\ll %\frac{y\rho(u)^{\eta(1-\nu)}}{\phi(k)}.$$
%\end{theorem}
%The $k=1$ case gives us an improvement over Hildebrand's result.
%\begin{varthmHilcor}
%Let $Q$, $y$, and $x$ be as in the previous theorem.  Then
%$$\sum_{\substack{x\leq n\leq x+y \\ n\in \mathcal S(Q)}}1\ll y\rho(u_x)^{C_1-(1-C_1)\eta+o(1)}.$$
%Equivalently,
%$$\sum_{\substack{x\leq n\leq x+y \\ n\in \mathcal S(Q)}}1\ll y\rho(u_y)^{1+C_1\left(\frac 1\eta-1\right)+o(1)}.$$
%\end{varthmHilcor}

\section{Applications}

We also present two applications of these theorems.  First, letting $\tau_d$ denote the $d$-fold divisor function (and letting $\tau=\tau_2$), we find a bound for the sum of $\tau_d(n)^R$ as $R$ slowly goes to infinity:

\begin{theorem}\label{dfold}
Let $x$, $y$, and $k$ be as in Theorem \ref{Maintheorem}.  Choose an integer $d\geq 2$, and let $R=\frac{\log\log\log x}{g(x)}$ for some $g(x)$ such that $g(x)\to\infty$ as $x\to\infty$.  Then for any $\epsilon>0$,
$$\sum_{\substack{x\leq n\leq x+y \\ n\equiv a\pmod k}}\tau_d(n)^{R}\ll \frac{y(\log x)^\epsilon}{\phi(k)}\left(\frac{\phi(k)}{k}\log x\right)^{d^{R}-1}.$$
Moreover,
$$\sum_{\substack{x\leq n\leq x+y \\ n\equiv a\pmod k \\ n\in \mathcal S(Q)}}\tau_d(n)^{R}\ll \rho(u)^{C_2+o(1)}\frac{y(\log x)^\epsilon}{\phi(k)}\left(\frac{\phi(k)}{k}\log x\right)^{d^{R}-1}.$$
\end{theorem}


For the second application, we note that one of the inspirations for this paper comes from a 2008 result of Soundararajan on smooth numbers \cite{So}.  We state this result here.

\begin{theorem}[Soundararajan 2008]\label{SoundThm}
Assume the Riemann Hypothesis. Let $x$ be large, and suppose that $Q\geq \exp\left(5\sqrt{\log x\log\log x}\right)$.  There is an absolute constant $B$ such that with $y = Bu\sqrt x/\rho(u/2)$ we have
\begin{gather}\label{Soundbound}
\psi(x+y,Q)-\psi(x,Q)\gg zx^{-\epsilon}.
\end{gather}
\end{theorem}
More recently, Goudout \cite{Go} found an unconditional version of this theorem, albeit with a smaller range for $Q$ and $z$ but a better lower bound:
\begin{theorem}[Goudout 2017]\label{GoThm}
Fix a positive $\epsilon<\frac 16$ and a constant $u_0$ depending on $\epsilon$. Let $Q$ be such that $\exp\left((\log x)^{\frac 56+\epsilon}\right)\leq Q\leq x^{\frac 1{u_0}}$.  Then for any $h$ with $\rho(u)^{-3-\epsilon}\leq h\leq \sqrt x$, if we let $y=h\sqrt x$ then
\begin{gather}\label{Gobound}
\psi(x+y,Q)-\psi(x,Q)\gg y\frac{\rho(u)^2}{\log^3 x}.
\end{gather}
\end{theorem}
With our improvement on Shiu's theorem, we can improve the lower bound in (\ref{Soundbound}) to almost the correct order of magnitude, and we can similarly improve the exponent on $\rho(u)$ in (\ref{Gobound}) to $1+o(1)$:
\begin{theorem}\label{smooththm}
Let $Q$, $x$, and $y$ satisfy either the hypotheses in either Theorem \ref{SoundThm} or \ref{GoThm}, where in the former case we also assume the Riemann Hypothesis.  Additionally, assume that $u\geq (\log\log x)^2.$  Then there exists a constant $C_3$ such that 
%Assume the Riemann Hypothesis, and let $x$, $Q$, $z$, and $B$ be as in Theorem \ref{SoundThm}.  
\begin{gather}
\psi(x+y,Q)-\psi(x,Q)\gg y\rho(u)^{1+\frac{C_3\log\log \log\log x}{\log\log \log x}}.
\end{gather}
In particular, we can take $C_3>1-C_2$.
\end{theorem}
The assumption here that $u\geq (\log\log x)^2$ is simply to eliminate log factors, since if $u$ satisfies this inequality then for any $A$, $$\log^A x=\rho(u)^{o(1)}.$$

We note that Jain \cite{Ja} has proven a lower bound for $Q$-smooth numbers in short intervals with $Q$ as small as $$\exp\left((\log x)^{\frac 23}(\log\log x)^{\frac 43+\epsilon}\right)$$ and $$y\geq \exp\left((1+\epsilon)\left(\frac{11}{16}u\log u+2\log\log x\right)\right).$$
However, the methods in that paper do not currently seem amenable to improvement by our methods, and hence we do not pursue this further.

\section{Outlines of proofs}
\subsection{Outline of Proof of Theorem \ref{Maintheorem}}

Shiu's original proof relies heavily on the following identity: for any $\delta>3/4$,
\begin{align*}
\sum_{\substack{n\leq x \\ (n,k)=1}}\frac{f(n)}{n^\delta}\ll &\prod_{\substack{p\leq x \\ p\nmid k}}\left(1+\frac{f(n)}{n^\delta}+\sum_{l=2}^\infty \frac{f(p^l)}{p^{l\delta}}\right)\leq \exp\left(\sum_{\substack{p\leq x \\ p\nmid k}}\left(\frac{f(p)}{p^\delta}+\sum_{l=2}^\infty \frac{f(p^l)}{p^{l\delta}}\right)\right).
\end{align*}
Since $f(p)$ is bounded by a constant $A_1$ in Shiu's work, we have that if $p^2>A_1$ then
\begin{gather}\label{psum}\sum_{l=2}^\infty \frac{f(p^l)}{p^{l\delta}}\leq \sum_{l=2}^\infty \left(\frac{A_1}{p^\delta}\right)^l\leq \frac{2A_1^2}{p^{2\delta}}\end{gather}
Meanwhile, if $p^2\leq A_1$ then (by Condition ii) we have $f(p^l)\leq A_2p^{l\epsilon}$, which means that 
\begin{gather}\label{psum2}\sum_{l=2}^\infty \frac{f(p^l)}{p^l}\leq \sum_{l=2}^\infty A_2\left(\frac{p^\epsilon}{p^\delta}\right)^l\leq \frac{3A_2}{p^{2\delta-2\epsilon}}.\end{gather}
Hence 
\begin{gather}\label{psum3}
\sum_p \sum_{l=2}^\infty \frac{f(p^l)}{p^{l\delta}}\ll 1,
\end{gather}
and thus
\begin{align}\label{psum3simple}
\sum_{\substack{n\leq x \\ (n,k)=1}}\frac{f(n)}{n^\delta}\ll \exp\left(\sum_{\substack{p\leq x \\ p\nmid k}}\frac{f(p)}{p^\delta}\right).
\end{align}
%For numbers $n$ with small numbers of prime factors, 
Shiu then splits the sum 
$$\sum_{\substack{x\leq n<x+y \\ n\equiv a\pmod k}}f(n)$$
into four subsums.  Three of the sums can be handled by sieve arguments, bounds on smooth numbers, or similar.  For the fourth sum, however, Shiu uses an alteration to (\ref{psum3simple}), finding that for relatively small $r$,
\begin{align}\label{Lemma3}
\sum_{\substack{n\geq z^\frac 12 \\ p(n)\leq z^\frac 1r \\ (n,k)=1}}\frac{f(n)}{n}\ll \exp\left(\sum_{\substack{p\leq x \\ p\nmid k}}\frac{f(p)}{p}-\frac{1}{10}r\log r\right).
\end{align}
This comes in handy when $n$ has a ``reasonable" number of distinct prime factors $r$, which is the set of $n$ that Shiu denoted by IV in his paper.  In this case, the negative term will counteract the growing number of factors.  In particular, for a constant $A_5$ that depends on both $A_1$ and the choice of $z$, and for an $r_0$ depending on $x$, he is able to bound the sum of $f(n)$ in this set with
$$\sum_{n\in IV}f(n)\ll \sum_{2\leq r\leq r_0}\frac{y}{\phi(k)\log x}A_5^r\exp\left(\sum_{\substack{p\leq x \\ p\nmid k}}\frac{f(p)}{p}-\frac{1}{10}r\log r\right)\ll \frac{y}{\phi(k)\log x}\exp\left(\sum_{\substack{p\leq x \\ p\nmid k}}\frac{f(p)}{p}\right),$$
which yields the main theorem in Shiu's paper.

If $f(p)$ is an increasing function of $x$, however, a few of these arguments fall apart.  For one, (\ref{psum3}) no longer holds, since $A_1$ and $A_2$ in (\ref{psum}) and (\ref{psum2}) are replaced by functions of $x$.  The bound in (\ref{Lemma3}) also no longer holds for much the same reason.  Moreover, in the bound for the sum over IV, $A_5^r$ is replaced by $\max_{p\leq x}\left(A_5 f(p)\right)^r$, which is a function on $r$ that may not be smaller than $\exp\left(\frac{1}{10}r\log r\right)$.  

In order to close these gaps in the proof, we make a couple of alterations.  A key change is that, for a parameter $J=(\log \log x)^{1-\beta}$, we split $n$ into $J$-smooth and $J$-rough parts, which we will call $s$ and $t$, respectively.   Since $f$ is multiplicative, we have $f(n)=f(s)f(t)$.  The sum over $s$ is small and relatively easy to handle, since $s$ itself is usually small.  For the sum over $t$, we now have
\begin{gather}\label{psum4}\sum_{p\geq J}\sum_{l=2}^\infty \frac{f(p^l)}{p^{l\delta}}\ll \sum_{p\geq J}\sum_{l=2}^\infty \left(\frac{A_1(\log\log x)^\beta}{p^\delta}\right)^l\leq 2\sum_{p\geq J}\frac{A_1^2(\log\log x)^{2\beta}}{p^{2\delta}}\ll \frac{A_1^2(\log\log x)^{2\beta}}{J^{2\delta-1}}=o(1),\end{gather}
since $\beta$ is much smaller than 1/5.  Hence,
%and thus
\begin{align}\label{psum4simple}
\sum_{\substack{n\leq x \\ (n,k)=1 \\ n\in \mathcal T}}\frac{f(n)}{n^\delta}\ll \exp\left(\sum_{\substack{J<p\leq x \\ p\nmid k}}\frac{f(p)}{p^\delta}\right).
\end{align}
In the cases where we have a larger number $r$ of distinct prime factors (i.e. the case where $n\in IV$), we will now have a slightly worsened bound of
\begin{gather}\label{Lemma3a}\sum_{n\in IV}f(n)\ll \frac{y}{\phi(k)\log x}(A_5 (\log\log x)^{\beta})^\frac{20r}{\alpha\kappa}\exp\left(\sum_{\substack{p\leq x \\ p\nmid k}}\frac{f(p)}{p}-\frac 12(1-\frac 54\beta)r\log r+2A_1(\log\log x)^\beta r^{1-\frac{5}{4}\beta}\right).
\end{gather}
The $\frac 12(1-\frac 54\beta)r\log r$ will eventually dominate once $r$ is large enough, but this will not occur until $r$ is close to $\log\log x$.  This transition period from $r=2$ to $r$ being close to $\log\log x$ is one of the reasons for the loss of $(\log x)^\epsilon$ in the main theorem.

\subsection{Outline of Proof of Theorem \ref{MnThm2}}

The case where $f$ is smooth-supported is a bit simpler, since in this case it is not possible for $r$ to be small.  After all, if $f(n)$ is only supported on $Q$-smooth numbers then $n$ must have at least $\frac{\log x}{\log Q}$ factors.  Hence we can use either (\ref{Lemma3}) if $f\in M_Q$ or (\ref{Lemma3a}) with $r\geq \frac{\log x}{\log Q}$ if $f\in M_Q(\beta)$.


\section{Removing Smooth Numbers}

Now, we begin the proof of Theorem \ref{Maintheorem}.   We first show that the contribution of numbers $n$ with a large smooth factor will be minimal.  

As above, let $$J=(\log \log x)^{1-\beta}.$$
We denote by $\mathcal S$ the $J$-smooth numbers and by $\mathcal T$ the $J$-rough numbers.  Our decomposition of $n$ is now $n=st$ where $s\in \mathcal S$ and $t\in \mathcal T$.  (In fact, we will generally use the convention that $s$ and $t$ indicate numbers in $\mathcal S$ and $\mathcal T$, respectively.)  Then
\begin{align*}
\sum_{\substack{x\leq n\leq x+y \\ n\equiv a\pmod k}}f(n)=&\sum_{s\leq (\log x)^{10}}\sum_{\substack{\frac xs\leq t\leq \frac xs+\frac ys \\ t\equiv a\overline{s}\pmod k}}f(st)+\sum_{(\log x)^{10}<s\leq x^\alpha}\sum_{\substack{\frac xs\leq t\leq \frac xs+\frac ys \\ t\equiv a\overline{s}\pmod k}}f(st)+\mathop{\sum\sum}_{\substack{x\leq st\leq x+y \\ st\equiv a\pmod k \\ s>x^\alpha}}f(st)\\
=&V_1+V_2+V_3,
\end{align*}
say.  

For any $z$ such that $J\leq (\log z)^{1-\beta}$, we use the identity that
\begin{gather}\label{psizJ}
\psi(z,J)=\left(\frac{\log z}{J}\right)^{(1+o(1))\pi(J)}=\left(\frac{\log z}{(\log\log x)^{1-\beta}}\right)^{(1+o(1))\frac{(\log\log x)^{1-\beta}}{(1-\beta)\log\log \log x}}\ll z^\epsilon
\end{gather}
for any $\epsilon>0$ \cite[(1.18)]{Gr}.

Using this, we can quickly handle $V_3$.
\begin{lemma}\label{LemmaV3}
$$V_3\ll \frac{y}{qx^{\frac{\alpha}{3}}}.$$
\end{lemma}
\begin{proof}
Note that, since $s$ is $J$-smooth, if $s>x^\alpha$ then $s$ can be split into $s=s_1s_2$, where $x^{\frac{\alpha}{2}}\leq s_1\leq x^{\frac{\alpha}{2}}J$.  (This decomposition is not necessarily unique, but this will not matter.)  Write $n'=s_2t$.  Recalling that $f(n)\ll n^\epsilon$ when $n>\log x$, we have
$$V_3\ll x^{\epsilon}\sum_{x^{\frac{\alpha}{2}}\leq s_1\leq x^{\frac{\alpha}{2}}J}\sum_{\substack{\frac x{s_1}\leq n'\leq \frac{x+y}{s_1} \\ n'\equiv a\overline{s_1}\pmod k}}1\ll x^{\epsilon}\sum_{x^{\frac{\alpha}{2}}\leq s_1\leq x^{\frac{\alpha}{2}}J}\frac{y}{qs_1}\ll \frac{y}{qx^{\frac{\alpha}{2}-\epsilon}}\sum_{x^{\frac{\alpha}{2}}\leq s_1\leq x^{\frac{\alpha}{2}}J}1.$$
By (\ref{psizJ}), we know that 
$$\psi(x^\frac \alpha 2 J,J)\ll x^\epsilon$$
for any $\epsilon>0$.

Thus,
$$V_3\ll \frac{y}{qx^{\frac{\alpha}{3}}}.$$
\end{proof}

In order to bound the other two sums, we will require the following theorem, which we will prove later.
\begin{theorem}\label{roughbd}
Let $f\in M(\beta)$.  Then for any $\epsilon>0$,
$$\sum_{\substack{x\leq t\leq x+y \\ t\equiv a\pmod k}}f(t)\ll \frac{y}{\phi(k)(\log x)^{1-\epsilon}}\left(\sum_{\substack{J<p\leq x \\ p\nmid k}}\frac{f(p)}{p}\right).$$
\end{theorem}

We also prove a variant for smooth-supported functions:
\begin{theorem}\label{roughbd2}
%Let $f\in M_Q$.  
%Then there exists $C_1$ such that
%\begin{gather}
%\sum_{\substack{x\leq n\leq x+y \\ n\equiv a\pmod k}}f(n)\ll \frac{y\rho(u)^{C_1+o(1)}}{\phi(k)(\log x)^{1-\epsilon}}\left(\sum_{\substack{J<p\leq x \\ p\nmid k}}\frac{f(p)}{p}\right).
%\end{gather}
%Moreover, 
Let $f\in M_Q(\beta)$. Then there exists a $C_2$ such that for any $\epsilon>0$,
\begin{gather}\label{rbdpart1}\sum_{\substack{x\leq t\leq x+y \\ t\equiv a\pmod k}}f(t)\ll \frac{y\rho(u)^{C_2+o(1)}}{\phi(k)(\log x)^{1-\epsilon}}\left(\sum_{\substack{J<p\leq x \\ p\nmid k}}\frac{f(p)}{p}\right).
\end{gather}
\end{theorem}
%Note that in the former case, we do not need to split $n$ into $st$, since $f(p)$ is $O(1)$.  Hence, the sum over $s\in \mathcal S$ is irrelevant there.

%Recalling that
%$$\psi(z,J)\ll X\left(\frac{\log z}{\log J}\right)^{-\frac{\log z}{\log J}},$$
%we see that if  $x^{\frac{\alpha}{2}}\leq z\leq x^{\frac{\alpha}{2}}J$ then
%$$\psi(z,J)\ll z\left(\frac{\alpha\log x}{2\delta\log \log x}\right)^{-\frac{\alpha\log x}{2\delta\log \log x}}= zx^{-\frac{\alpha}{2\delta}(1+o(1))}.$$
%So by partial summation,
%\begin{align*}\sum_{x^{\frac{\alpha}{2}}\leq s_1\leq x^{\frac{\alpha}{2}}J}\frac{1}{s_1}=&\sum_{x^{\frac{\alpha}{2}}\leq t\leq x^{\frac{\alpha}{2}}J}\left(\frac{1}{t}-\frac{1}{t+1}\right)\sum_{x^{\frac{\alpha}{2}}\leq s'\leq t}1 +\frac{1}{x^{\frac{\alpha}{2}}J}\sum_{x^{\frac{\alpha}{2}}\leq s_1\leq x^{\frac{\alpha}{2}}J}1\ll x^{-\frac{\alpha}{3\delta}},
%\end{align*}
%where $t$ indicates that we have dropped the restriction of being in $\mathcal S$.  Taking $\epsilon<\frac{\alpha}{12\delta}$ in (\ref{Condition2}), we then have
%$$V_3\ll \frac{yx^{-\frac{\alpha}{4\delta}}}{q}.$$
%


If we are willing to assume these, we can easily find bounds for $V_1$ and $V_2$.  For ease of notation, define $\varrho$ such that $\varrho=1$ if $f\in M(\beta)$ and $\varrho=\rho(u)^{C_2+o(1)}$ if $f\in M_Q(\beta)$.

\begin{lemma}
Assume Theorem \ref{roughbd} or \ref{roughbd2}.  Then
$$V_1\ll \frac{y\varrho}{\phi(k)(\log x)^{1-\epsilon_0}}\exp\left(\sum_{\substack{J<p\leq x \\ p\nmid k}}\frac{f(p)}{p}\right).$$
\end{lemma}
\begin{proof}
Clearly, 
$$\sum_{s\leq (\log x)^{10}}\frac 1s\ll \log\log x.$$
So applying either theorem and Condition \ref{Condition2} gives
\begin{align*}
\sum_{s\leq (\log x)^{10}}\sum_{\substack{\frac xs\leq t\leq \frac xs+\frac ys \\ t\equiv a\overline{s}\pmod k}}f(st)\ll & \sum_{s\leq (\log x)^{10}}\frac{y\varrho f(s)}{\phi(k)s(\log x)^{1-\epsilon}}\exp\left(\sum_{\substack{J<p\leq x \\ p\nmid k}}\frac{f(p)}{p}\right)\\
\ll & \frac{y\varrho }{\phi(k)(\log x)^{1-\epsilon_0}}\exp\left(\sum_{\substack{J<p\leq x \\ p\nmid k}}\frac{f(p)}{p}\right),
\end{align*}
which is as required.
\end{proof}
Finally, we handle $V_2$.
\begin{lemma}
Assume Theorem \ref{roughbd} or \ref{roughbd2}.  Then for any $\epsilon_0>0$,
$$V_2\ll \frac{y\varrho }{\phi(k)(\log x)^{10-\epsilon_0}}\exp\left(\sum_{\substack{J<p\leq x \\ p\nmid k}}\frac{f(p)}{p}\right).$$
\end{lemma}
\begin{proof}
We begin with the assumed theorem:
\begin{gather}\label{T2sum}\sum_{(\log x)^{10}<s\leq x^\frac \alpha 2}\sum_{\substack{\frac xs\leq t\leq \frac xs+\frac ys \\ t\equiv a\overline{s}\pmod k}}f(st)\ll \sum_{(\log x)^{10}<s\leq x^\alpha}\frac{y\varrho f(s)}{\phi(k)s(\log x)^{1-\epsilon}}\exp\left(\sum_{\substack{J<p\leq x \\ p\nmid k}}\frac{f(p)}{p}\right).\end{gather}
Note that for any $s$ in the interval $((\log x)^{10},x^\alpha]$, we have 
$$(\log\log x)^{1-\beta}\leq \left(\frac{1}{10}\log s\right)^{1-\beta}.$$
So by (\ref{psizJ}),
\begin{gather}\label{zpsi}
\psi(s,J)\ll s^{\epsilon}
\end{gather}
for any $\epsilon>0$.

So we apply partial summation to the sum $V_2$, where we will use the variable $v$ to indicate where we have dropped the requirement of smoothness:
\begin{align*}\sum_{(\log x)^{10}<s\leq x^\frac \alpha 2}\frac{f(s)}{s}=&\sum_{(\log x)^{10}<v\leq x^\frac \alpha 2-1}\left(\frac{f(v)}{v}-\frac{f(v+1)}{v+1}\right)\sum_{(\log x)^{10}<s\leq v}1 +\frac{1}{x^{\frac \alpha 2}}\sum_{(\log x)^{10}<s\leq x^\frac \alpha 2}f(s).
\end{align*}
Applying (\ref{zpsi}) to the sum and $f(v)\ll v^\epsilon$ to the function then gives
\begin{align*}\sum_{(\log x)^{10}<s\leq x^\frac \alpha 2}\frac{f(s)}{s}\ll &\left(\sum_{(\log x)^{10}<v\leq x^\frac \alpha 2-1}\frac{1}{v^{1-2\epsilon}(v+1)} \right)+\frac{1}{x^{\frac \alpha 2-2\epsilon}}
\ll (\log x)^{-9}.
\end{align*}
Plugging this into (\ref{T2sum}) then yields the lemma.

\end{proof}

The next sections will then set up the machinery to prove Theorems \ref{roughbd} and \ref{roughbd2}.
\section{The Sum over Rough Numbers}

Hereafter, for ease of notation, we will write $X=x/s$, $Y=y/s$, and $b\equiv a\overline{s}\pmod k$.  We are tasked with evaluating
$$\sum_{\substack{X\leq t\leq X+Y \\ t\equiv b\pmod k}}f(t).$$

From here, our methods will largely follow those of \cite{Sh}.  As in that paper, we will let $p(n)$ and $q(n)$ denote the greatest and least prime factors of $n$, respectively, and we define
$$\Phi(x,y,z;k,a)=\sum_{\substack{x\leq n\leq x+y \\ n\equiv a \pmod k \\ q(n)>z}}1.$$
By Lemma 2 of \cite{Sh}, we know that 
\begin{gather}\label{rough}\Phi(x,y,z;k,a)\ll \frac{y}{\phi(k)\log z}+z^2.\end{gather}
Generally, we will be taking $y$, $k$, and $z$ such that $y/k$ is much larger than $z$, and hence the latter term can be ignored.

We also record a few statements about smooth numbers for later.  From \cite[Lemma 1]{Sh},
\begin{gather}\label{psibound}
\psi(x,\log x\log\log x)\ll \exp\left(\frac{3\log x}{(\log \log x)^\frac 12}\right).
\end{gather}
Meanwhile, for larger $z$, we know from (\ref{Di}) that
\begin{gather}\label{smoothbd}
\psi(x,z)\ll x\exp\left(-\frac{\log x(\log x-\log z)}{\log z}\right).
\end{gather}
%We will use repeatedly the fact that
%\begin{gather}\label{plsum}\sum_{l=2}^\infty\sum_{J<p\leq x}\frac{f(p^l)}{p^l}\ll \sum_{l=2}^\infty\sum_{J<p\leq x}\frac{(\log\log x)^{l\beta}}{p^l}\ll \frac{1}{(\log\log x)^{1-2\beta}},\end{gather}
%by Condition 1 and our choice of $J$.
%Since
%\begin{align*}
%\sum_{\substack{t\leq X \\ (t,k)=1}}\frac{f(t)}{t}\ll &\prod_{\substack{J<p\leq x \\ p\nmid k}}\left(1+\frac{f(p)}{p}+\sum_{l=2}^\infty \frac{f(p^l)}{p^l}\right)\leq \exp\left(\sum_{\substack{J<p\leq x \\ p\nmid k}}\left(\frac{f(p)}{p}+\sum_{l=2}^\infty \frac{f(p^l)}{p^l}\right)\right),
%\end{align*}
%we can apply (\ref{plsum}) to find
%\begin{align}\label{expto}
%\sum_{\substack{t\leq X \\ (t,k)=1}}\frac{f(t)}{t}\ll \exp\left(\sum_{\substack{J<p\leq x \\ p\nmid k}}\frac{f(p)}{p}\right).
%\end{align}
We will also require one more preliminary lemma before we turn our attention to the main sum.  This is a slightly altered version of \cite[Lemma 4]{Sh}.

\begin{lemma}\label{lemmalogr}
Let $f\in \mathcal M(\beta)$.  Then as $z\to\infty$,
$$\sum_{\substack{n\geq z^\frac 12 \\ p(n)\leq z^\frac 1r \\ (n,k)=1 \\ n\in \mathcal T}}\frac{f(n)}{n}\ll \exp\left(-\frac 12\left(1-\frac{5}{4}\beta\right)r\log r+\sum_{\substack{J<p\leq z^\frac 1r \\ p\nmid k}}\frac{f(p)}{p}+2A_1 (\log\log x)^{\beta}r^{1-\frac{5}{4}\beta}\right)$$
uniformly in $k$ and $r$, provided that $1\leq r\leq \frac{\log z}{4\log\log z}$.
\end{lemma}
\begin{proof}
Let $\frac 34 <\delta\leq 1$.  By Rankin's trick,
\begin{align*}
\sum_{\substack{n\geq w \\ p(n)\leq v \\ (n,k)=1 \\ n\in \mathcal T}}\frac{f(n)}{n}\leq &w^{\delta-1}\sum_{\substack{n\geq w \\ p(n)\leq v \\ (n,k)=1 \\ n\in \mathcal T}}\frac{f(n)}{n^\delta}\leq w^{\delta-1}\sum_{\substack{n\geq 1 \\ p(n)\leq v \\ (n,k)=1 \\ n\in \mathcal T}}\frac{f(n)}{n^\delta}\\
=&w^{\delta-1}\prod_{\substack{J<p\leq v \\ p\nmid k}}\left(1+\frac{f(p)}{p^\delta}+\sum_{l=2}^\infty \frac{f(p^l)}{p^{l\delta}}\right)\ll \exp\left((\delta-1)\log w+\sum_{\substack{J<p\leq v \\ p\nmid k}}\frac{f(p)}{p^\delta}\right)
\end{align*}
by (\ref{psum4simple}).
We rewrite $\frac{f(p)}{p^\delta}$ as $\frac{f(p)}{p}+\frac{f(p)}{p}(p^{1-\delta}-1)$ and note that
\begin{align*}
\sum_{\substack{J<p\leq v \\ p\nmid k}}\frac{f(p)}{p}(p^{1-\delta}-1)\leq &A_1 (\log\log x)^{\beta}\sum_{p\leq v}\frac 1p\sum_{n=1}^\infty \frac{\left((1-\delta)(\log p)\right)^n}{n!}\\
\leq &A_1 (1-\delta)(\log\log x)^{\beta}\sum_{n=1}^\infty \frac{\left((1-\delta)(\log v)\right)^{n-1}}{n!}\sum_{p\leq v}\frac{\log p}p\\
\leq &A_1 (\log\log x)^{\beta}\sum_{n=1}^\infty \frac{\left((1-\delta)(\log v)\right)^{n}}{n!}\\
\leq &2A_1 (\log\log x)^{\beta}v^{1-\delta}.
\end{align*}
So
\begin{align*}
\sum_{\substack{n\geq u \\ p(n)\leq v \\ (n,k)=1 \\ n\in \mathcal T}}\frac{f(n)}{n}\ll \exp\left((\delta-1)\log u+\sum_{\substack{J<p\leq v \\ p\nmid k}}\frac{f(p)}{p}+2A_1 (\log\log x)^{\beta}v^{1-\delta}\right).
\end{align*}



Letting $u=z^\frac 12$, $v=z^\frac 1r$, and $\delta=1-\frac{(1-\frac 54\beta )r\log r}{\log z}$ then gives

\begin{align*}
\sum_{\substack{n\geq u \\ p(n)\leq v \\ (n,k)=1 \\ n\in \mathcal T}}\frac{f(n)}{n}\ll \exp\left(-\frac 12\left(1-\frac 54\beta \right)r\log r+\sum_{\substack{J<p\leq z^\frac 1r \\ p\nmid k}}\frac{f(p)}{p}+2A_1 (\log\log x)^{\beta}r^{1-\frac 54\beta }\right).
\end{align*}
\end{proof}

\section{Partitioning the Sum: The Proof of Theorem \ref{roughbd}}
Now, we partition the sum as in Section 5 of \cite{Sh}.  Let 
$$z=Y^{\frac{\alpha}{10}}.$$
Note that by this definition, $Y/k>z^2$, and hence we can omit the $z^2$-term in (\ref{rough}).

For each $n$ with $X\leq n<X+Y$, $n\equiv b\pmod k$, write
$$n=p_1^{u_1} \cdots  p_j^{u_j}p_{j+1}^{u_{j+1}}\cdots p_l^{u_l}=c_nd_n$$
where $p_1<p_2<\cdots <p_l$ and $c_n$ is the product of the first $j$ terms, where $j$ is chosen such that
$$c_n\leq z<c_np_{j+1}^{u_{j+1}}.$$
We then define the following four subclasses of $\mathcal T$:
\begin{gather*}
\mbox{I}:q(d_n)>z^\frac 12,\\
\mbox{II}:q(d_n)\leq z^\frac 12,c_n\leq z^\frac 12,\\
\mbox{III}:q(d_n)\leq \left(\log x\log \log x\right)^5,c_n>z^\frac 12,\\
\mbox{IV}:\left(\log x\log \log x\right)^5<q(d_n)\leq z^\frac 12,c_n>z^\frac 12.
\end{gather*}
\subsection{Class I}
Here, we have
$$\sum_{n\in I}f(n)\leq \sum_{\substack{c\leq z \\ c\in \mathcal T
}}f(c)\sum_{\substack{\frac Xc\leq d\leq \frac{X+Y}{c} \\ d\equiv b\overline{c}\pmod k \\ q(d)>z^\frac 12}}f(d).$$
Note that for any such $n$, $q(d)>z^\frac 12=Y^{\frac{\alpha}{20}}=X^{\frac{\alpha\kappa}{20}}.$  So $d$ can have at most $\frac{20}{\alpha\kappa}$ prime factors, and hence $f(d)\leq \left(A_1(\log\log x)^\beta\right)^{\frac{20}{\alpha\kappa}}$.  Noting that the number of $z^\frac 12$-rough numbers up to $x$ in a congruence class mod $k$ is $\ll \frac{y}{\phi(k)\log z}$, and recalling that $\beta<\frac{\alpha\kappa}{41}$, we can bound the above with
\begin{align*}
\sum_{n\in I}f(n)\leq &\left(A_1(\log\log x)^\beta\right)^{\frac{20}{\alpha\kappa}}\sum_{\substack{c\leq z \\ c\in \mathcal T}}f(c)\sum_{\substack{\frac Xc\leq d\leq \frac{X+Y}{c} \\ d\equiv b\overline{c}\pmod k \\ q(d)>z^\frac 12}}1\\
\ll &\frac{Y}{\phi(k)\log z}\left(\log\log x\right)^{\frac{20}{41}}\sum_{\substack{c\leq z \\ c\in \mathcal T}}\frac{f(c)}{c}\\
\ll &\frac{Y}{\phi(k)\log x}\left(\log\log x\right)^{\frac{20}{41}}\exp\left(\sum_{\substack{J<p\leq X \\ p\nmid k}}\frac{f(p)}{p}\right),
\end{align*}
where the last line is by (\ref{psum4simple}).  

%This can be simplified to 
%$$\ll \frac{Y}{\phi(k)(\log x)^{1-\epsilon}}\exp\left(\sum_{\substack{J<p\leq X \\ p\nmid k}}\frac{f(p)}{p}\right)$$
%for every $\epsilon>0$.

\subsection{Class II}
In this case, there must exist a $p|n$ such that $p\leq z^\frac 12$ and $p^u>z^\frac 12$.  For such a $p$, let $u_p$ denote the least integer such that $p^{u_p}>z^\frac 12$.  Then $p^{-u_p}\leq \min\{z^{-\frac 12},p^{-2}\}$.  Hence
$$\sum_{p\leq z^\frac 12}\frac{1}{p^{u_p}}\leq \sum_{p\leq z^\frac 14}z^{-\frac 12}+\sum_{p>z^\frac 14}\frac 1{p^2}\ll z^{-\frac 14}.$$
Then for any $\epsilon>0$,
$$\sum_{n\in II}f(n)\ll x^\epsilon\sum_{p\leq z^\frac 12}\sum_{\substack{X\leq n<X+Y \\ n\equiv b\pmod k \\ p^{u_p}|n }}1\ll x^\epsilon\sum_{p\leq z^\frac 12}\frac{Y}{kp^{u_p}}\ll \frac{Yx^\epsilon}{kz^{\frac 14}}.$$
\subsection{Class III}
In this case, we have $c$ such that $z^\frac 12<c\leq z$ and $c$ is $(\log x\log \log x)^5$-smooth.  So 
$$\sum_{n\in III}f(n)\leq \sum_{\substack{z^\frac 12<c\leq z \\ p(c)\leq (\log x\log \log x)^5}}f(c)\sum_{\substack{\frac Xc\leq d\leq \frac{X+Y}{c} \\ d\equiv b\overline{c}\pmod k }}f(d)\ll \sum_{\substack{z^\frac 12<c\leq z \\ p(c)\leq (\log x\log \log x)^5}}\frac{Yx^\epsilon}{kc}$$
Let $$v=\left(\frac{\log z^\frac 12}{5\log\left(\log x\log\log x\right)}\right)^{-\frac{\log z^\frac 12}{5\log\left(\log x\log\log x\right)}}\ll \left(\log z\right)^{-\frac{\log z}{11\log\log x}}\ll x^{-\frac{\alpha\kappa}{250}}.$$
We can handle the sum dyadically and apply (\ref{smoothbd}), finding
\begin{gather}\label{dyad}\sum_{\substack{z^\frac 12<c\leq z \\ p(c)\leq (\log x\log \log x)^5}}\frac{1}{c}\leq \sum_{j=0}^{\frac{\log z}{\log 2}}\frac{1}{2^jz^\frac 12}\sum_{\substack{2^jz^\frac 12<c\leq 2^{j+1}z^\frac 12 \\ p(c)\leq (\log x\log \log x)^5}}1 \ll \sum_{j=0}^{\frac{\log z}{\log 2}}\frac{1}{2^jz^\frac 12}\left[2^{j+1}z^\frac 12v\right]\ll v\log z\ll x^{-\frac{\alpha\kappa}{300}}.
\end{gather}
So
$$\sum_{n\in III}f(n)\ll \frac{Yx^{\epsilon-\frac{\alpha\kappa}{300}}}{k}\ll \frac{Yx^{-\epsilon'}}{k}$$
for some $\epsilon'>0$.

\subsection{Class IV}
This class gives us
$$\sum_{n\in IV}f(n)\leq \sum_{\substack{z^\frac 12<c\leq z \\ c\in \mathcal T
}}f(c)\sum_{\substack{\frac Xc\leq d\leq \frac{X+Y}{c} \\ d\equiv b\overline{c}\pmod k \\ q(d)>p(c) \\ (\log x\log\log x)^5< q(d)\leq z^\frac 12}}f(d).$$
Let 
$$r_0=\left[\frac{\log z}{5\log(\log x\log \log x)}\right],$$
where we now have $\log x\log \log x>z^{\frac{1}{r_0+1}}$.  Let $r$ be such that $2\leq r\leq r_0$, and consider those $n$ for which $z^{\frac{1}{r+1}}<q(d_n)\leq z^\frac 1r$. For such $n$, we know that the number of prime factors of $d_n$ is bounded by
$$\Omega(d_n)\leq \frac{\log x}{\log(q(d_n))}\leq \frac{20r}{\alpha\kappa}.$$
So
$$f(d_n)\leq \left(A_1(\log\log x)^\beta\right)^{\frac{20r}{\alpha\kappa}}\leq  A_1^{\frac{20r}{\alpha\kappa}}\left(\log\log x\right)^{\frac{20r}{41}}.$$
Hence
\begin{align*}
\sum_{n\in IV}f(n)\leq  &\sum_{2\leq r\leq r_0}A_1^{\frac{20r}{\alpha\kappa}}\left(\log\log x\right)^{\frac{20r}{41}}\sum_{\substack{z^\frac 12<c\leq z \\ p(c)<z^\frac 1r \\ c\in \mathcal T}}f(c)\Psi\left(\frac Xc,\frac Yc,z^\frac{1}{r+1};k,b\overline{c}\right)\\
\leq  &\sum_{2\leq r\leq r_0}A_1^{\frac{20r}{\alpha\kappa}}\left(\log\log x\right)^{\frac{20r}{41}}\left(\frac{Y(r+1)}{\phi(k)\log z}\right)\sum_{\substack{z^\frac 12<c\leq z \\ p(c)<z^\frac 1r \\ c\in \mathcal T}}\frac{f(c)}{c}
\end{align*}
by (\ref{rough}).

Let $$r_1=(\log\log x)^{1-\frac{\beta}{2}}.$$

We can split the sum over $r$ into
$$\sum_{2\leq r\leq r_1-1}+\sum_{r_1\leq r\leq r_0}.$$

For the former sum over $r$, we bound the sum over $c$ using (\ref{psum4simple}):
\begin{align*}
\sum_{2\leq r\leq r_1-1}&A_1^{\frac{20r}{\alpha\kappa}}\left(\log\log x\right)^{\frac{20r}{41}}\left(\frac{Y(r+1)}{\phi(k)\log z}\right)\sum_{\substack{z^\frac 12<c\leq z \\ p(c)<z^\frac 1r \\ c\in \mathcal T}}\frac{f(c)}{c}\\
\ll &\sum_{2\leq r\leq r_1-1}\left(\frac{Y(r+1)}{\phi(k)\log z}\right)A_1^{\frac{20r}{\alpha\kappa}}\left(\log\log x\right)^{\frac{20r}{41}}\exp\left(\sum_{\substack{J<p\leq z \\ p\nmid k}}\frac{f(p)}{p}\right)\\
\ll &\left(\frac{Yr_1}{\phi(k)\log z}\right)A_1^{\frac{20r_1}{\alpha\kappa}}\left(\log\log x\right)^{\frac{20r_1}{41}}\exp\left(\sum_{\substack{J<p\leq z \\ p\nmid k}}\frac{f(p)}{p}\right)\\
\ll &\left(\frac{Y}{\phi(k)\log z}\right)\exp\left((\log\log x)^{1-\frac{\beta}{3}}+\sum_{\substack{J<p\leq z \\ p\nmid k}}\frac{f(p)}{p}\right)
\end{align*}
Clearly, this is $\ll \frac{Y}{\phi(k)(\log x)^{1-\epsilon}}$ for any $\epsilon>0$.

%For the remaining sum over $r$, we note that $c$ must be relatively smooth.  Proceeding as in our proof of SOMETHING, we have

Turning now to the interval where $r\geq r_1$, we can apply Lemma \ref{lemmalogr} to bound the remaining sum as
\begin{align*}
&\ll \sum_{r_1\leq r\leq r_0}A_1^{\frac{20r}{\alpha\kappa}}\left(\log\log x\right)^{\frac{20r}{41}}\left(\frac{Y(r+1)}{\phi(k)\log z}\right)\\
&\cdot \exp\left(-\frac 12\left(1-\frac 54\beta\right)r\log r+\sum_{\substack{J<p\leq z^\frac 1r \\ p\nmid k}}\frac{f(p)}{p}+2A_1 (\log\log x)^{\beta}r^{1-\frac 54\beta}\right)\\
&\ll  \sum_{r_1\leq r\leq r_0}\left(\frac{Y}{\phi(k)\log z}\right)\\
&\cdot \exp\left(\sum_{\substack{J<p\leq X \\ p\nmid k}}\frac{f(p)}{p}+\frac{20r\log A_1}{\alpha\kappa}+\frac{20r}{41}\log\log\log x+2\log r-\frac 12\left(1-\frac 54\beta \right)r\log r+2A_1 (\log\log x)^{\beta}r^{1-\frac 54\beta }\right).
\end{align*}
Let
\begin{gather*}g(r)=\frac 12\left(1-\frac 54\beta \right)\log r,\\
h(r)=\frac{20\log A_1}{\alpha\kappa}+\frac{20}{41}\log\log\log x+\frac{2\log r}{r}+2A_1 (\log\log x)^{\beta}r^{-\frac 54\beta }
%2A_1 (\log\log x)^{\beta}r^{-\frac 54\beta}+\frac{20}{41}\log\log\log x+\frac{20\log A_1}{\alpha\kappa}+\frac{2\log r}{r}.
\end{gather*}
So the above is
$$\ll \sum_{r_1\leq r\leq r_0}\left(\frac{Y}{\phi(k)\log z}\right)\exp\left(-r(g(r)-h(r))+\sum_{\substack{J<p\leq X \\ p\nmid k}}\frac{f(p)}{p}\right).$$
Now, since $\beta<\frac{\alpha\kappa}{41}\leq \frac{1}{164}$,
we see that
\begin{align*}g(r_1)=&\frac 12\left(1-\frac 74\beta+\frac 58\beta^2+o(1)\right)\log\log\log x\\
\geq &\left(\frac 12-\frac{1}{164}+o(1)\right)\log\log\log x.
\end{align*}
Meanwhile,
\begin{align*}h(r_1)=&\frac{20}{41}\log\log\log x+2A_1 (\log\log x)^{\beta}(\log \log x)^{-\frac 54\beta+\frac 58\beta^2}+O(1)\\
=&\left(\frac 12-\frac{1}{41}+o(1)\right)\log\log\log x.
\end{align*}
So 
$$g(r_1)-h(r_1)\geq \left(\frac{1}{60}+o(1)\right)\log\log\log x.$$
Noting also that $g'(r)>h'(r)$ for any $r\in [r_1,r_0]$, we can bound the above sum over $r$ as
\begin{align*}\ll &\sum_{r_1\leq r\leq r_0}\left(\frac{Y}{\phi(k)\log z}\right)\exp\left(-\frac{r}{60}\log\log\log x+\sum_{\substack{J<p\leq X \\ p\nmid k}}\frac{f(p)}{p}\right)\\
\ll &\left(\frac{Y}{\phi(k)\log z}\right)\exp\left(\sum_{\substack{J<p\leq X \\ p\nmid k}}\frac{f(p)}{p}\right)
\end{align*}
Putting these together,
\begin{align*}\sum_{n\in IV}f(n)\ll & \left(\frac{Y}{\phi(k)\log x}\right)\exp\left((\log\log x)^{1-\frac{\beta}{3}}+\sum_{\substack{J<p\leq X \\ p\nmid k}}\frac{f(p)}{p}\right)\\
\ll & \left(\frac{Y}{\phi(k)(\log x)^{1-\epsilon}}\right)\exp\left(\sum_{\substack{J<p\leq X \\ p\nmid k}}\frac{f(p)}{p}\right)
\end{align*}
for any $\epsilon>0$.  This completes the proof of Theorem \ref{roughbd} and hence of Theorem \ref{Maintheorem}.



\section{Alterations for Smooth-Supported Functions}

Here, we alter the steps of the previous section to prove the analogous theorems for smooth-supported functions.  We note first that $\frac{\log Q}{\log x}=o(1)$, else the two bounds in Theorem \ref{MnThm2} are just Shiu's theorem and Theorem \ref{Maintheorem}, respectively.

Class I cannot happen, since if $q(d_n)>z^\frac 12>Q$ then $n$ is not $Q$-smooth.  For Classes II and III, we can use the same bounds as in the non-smooth case.  So we require only class IV, as the first three classes give
$$\sum_{\substack{X\leq t\leq X+Y \\ t\equiv a\pmod k}}f(t)=\sum_{n\in IV}f(n)+O\left(\frac{Y}{kz^{\frac 14-\epsilon}}\right)$$

We consider first the case when $f\in M_Q$.  Since we no longer have an unbounded $f(p)$, we could apply Lemma 4 of Shiu's original paper \cite{Sh}, which states that if $f\in M$ then
\begin{gather}\label{Shlem}\sum_{\substack{n\geq z^\frac 12 \\ p(n)\leq z^\frac 1r \\ (n,k)=1 \\ n\in \mathcal T}}\frac{f(n)}{n}\ll \exp\left(-\frac{1}{10}r\log r+\sum_{\substack{p\leq z \\ p\nmid k}}\frac{f(p)}{p}\right).
\end{gather}
The proof there is very similar to our proof of Lemma \ref{lemmalogr} except that Shiu takes $\delta=1-\frac{r\log r}{4\log z}$, whereas we take $\delta=1-\frac{(1-\frac 54\beta )r\log r}{\log z}$.  If we were to apply our $\delta$ in Shiu's original Lemma 4, we would instead have
\begin{gather}\label{Shlem4}\sum_{\substack{n\geq z^\frac 12 \\ p(n)\leq z^\frac 1r \\ (n,k)=1 \\ n\in \mathcal T}}\frac{f(n)}{n}\ll \exp\left(-\frac 12\left(1-\frac 54\beta\right)r\log r+\sum_{\substack{p\leq z \\ p\nmid k}}\frac{f(p)}{p}\right),
\end{gather}
where the $\ll$ depends on our choice of $\beta$.

Now, as we are dealing with Class IV, we note that since $z^\frac 12<c\leq z$, we must have $d>x^{1-\frac{\alpha \kappa}{20}}$.  So if $n$ is such that $f(n)\neq 0$ then $n$ must have at least $\frac{\left(1-\frac{\alpha \kappa}{20}\right)\log x}{\log Q}$ prime factors.  Recalling that
$$\Omega(d_n)\leq \frac{20r}{\alpha \kappa},$$
we must then have
$$\frac{\alpha\kappa \left(1-\frac{\alpha \kappa}{20}\right)\log x}{20\log Q}\leq r$$
or else $n$ is not $Q$-smooth.  Let 
$$r_2=\frac{\alpha\kappa \left(1-\frac{\alpha \kappa}{20}\right)\log x}{20\log Q}=\frac{\alpha\kappa \left(1-\frac{\alpha \kappa}{20}\right)}{20}u.$$
Then 
\begin{align*}
\sum_{n\in IV}f(n)\leq  &\sum_{r_2\leq r\leq r_0}A_1^{\frac{20r}{\alpha\kappa}}\left(\frac{Y(r+1)}{\phi(k)\log z}\right)\sum_{\substack{z^\frac 12<c\leq z \\ p(c)<z^\frac 1r }}\frac{f(c)}{c}\\
\ll  &\sum_{r_2\leq r\leq r_0}\left(\frac{Y}{\phi(k)\log z}\right)\exp\left(-\frac 12\left(1-\frac 54\beta\right)r\log r+\sum_{\substack{p\leq Q \\ p\nmid k}}\frac{f(p)}{p}+\frac{20r}{\alpha\kappa}\log A_1+2\log r\right)\\
\ll  &\sum_{r_2\leq r\leq r_0}\left(\frac{Y}{\phi(k)\log z}\right)\exp\left(-\frac 12\left(1-\frac 32\beta\right)r\log r\right)\exp\left(\sum_{\substack{p\leq Q \\ p\nmid k}}\frac{f(p)}{p}\right)\\
\ll  &\left(\frac{Y}{\phi(k)\log z}\right)\exp\left(-\left(\frac 12-\beta\right)r_2\log r_2\right)\exp\left(\sum_{\substack{p\leq Q \\ p\nmid k}}\frac{f(p)}{p}\right).
\end{align*}
Noting that 
$$\exp\left(-\left(\frac 12-\beta\right)r_2\log r_2\right)=u^{-\frac{\alpha\kappa}{40}\left(1-2\beta\right)u},$$
we can let $$C_1=\frac{\alpha\kappa}{41}.$$ 
This then gives us 
$$\sum_{n\in IV}f(n)\ll \rho(u)^{C_1+o(1)}\left(\frac{Y}{\phi(k)\log x}\right)\left(\sum_{\substack{J<p\leq Q \\ p\nmid k}}\frac{f(p)}{p}\right),$$
which proves the first half of the theorem.

For the case of $f\in \mathcal M_Q(\beta)$, we proceed similarly, except that we use Lemma \ref{lemmalogr}.  So again,


\begin{align*}
\sum_{n\in IV}f(n)
&\leq  \sum_{r_2\leq r\leq r_0}A_1^{\frac{20r}{\alpha\kappa}}\left(\log\log x\right)^{\frac{20r}{41}}\left(\frac{Y(r+1)}{\phi(k)\log z}\right)\sum_{\substack{z^\frac 12<c\leq z \\ p(c)<z^\frac 1r \\ c\in \mathcal T}}\frac{f(c)}{c}\\
&\ll  \sum_{r_2\leq r\leq r_0}\left(\frac{Y}{\phi(k)\log z}\right)\exp\left(-r(g(r)-h(r))+\sum_{\substack{J<p\leq Q \\ p\nmid k}}\frac{f(p)}{p}\right),
\end{align*}
where $g$ and $h$ are as in the previous section.

Since $r_2>r_1$, we note as before that $g'(r)>h'(r)$ for all $r\in [r_2,r_0]$.  Here, however, we also have
\begin{gather*}g(r_2)=\frac 12\left(1-\frac 54\beta \right)\left[\log u+\log\left(\frac{\alpha\kappa \left(1-\frac{\alpha \kappa}{20}\right)}{20}\right)\right].
%2A_1 (\log\log x)^{\beta}r^{-\frac 54\beta}+\frac{20}{41}\log\log\log x+\frac{20\log A_1}{\alpha\kappa}+\frac{2\log r}{r}.
\end{gather*}
%$$g(r_2)=\frac 12\left(1-\frac 54\beta \right)\frac{\alpha\kappa \left(1-\frac{\alpha \kappa}{20}\right)}{20}u\left[\log\left(u\right)-\log\left(\frac{\alpha\kappa \left(1-\frac{\alpha \kappa}{20}\right)}{20}\right)\right],$$
Moreover, since $\frac{\log x}{\log Q}\geq \log \log x$ by the assumption in Theorem \ref{MnThm2}, we know that
%$$h(r)=\frac{20\log A_1}{\alpha\kappa}+\frac{20}{41}\log\log\log x+\frac{2\log r}{r}+2A_1 (\log\log x)^{\beta}r^{-\frac 54\beta }$$
\begin{align*}
h(r_2)=&\frac{20}{41}\log\log\log x+O(1)+2A_1 (\log\log x)^{\beta}\left(\frac{\alpha\kappa \left(1-\frac{\alpha \kappa}{20}\right)}{20}u\right)^{-\frac 54\beta }\\
=&\left(\frac{20}{41}+o(1)\right)\log\log\log x\leq \left(\frac{20}{41}+o(1)\right)\log u.\end{align*}
%\begin{align*}
%h(r_2)\leq &2A_1 \left(\frac{\alpha\kappa \left(1-\frac{\alpha \kappa}{20}\right)}{20}\frac{\log x}{\log Q}\right)^{1-\frac \beta 4}+\frac{20}{41}\left(\frac{\alpha\kappa \left(1-\frac{\alpha \kappa}{20}\right)}{20}\frac{\log x}{\log Q}\right)\log \left(\frac{\log x}{\log Q}\right)+2\log\log x\\
%= &\left(\frac{20}{41}+o(1)\right)\left(\frac{\alpha\kappa \left(1-\frac{\alpha \kappa}{20}\right)}{20}\frac{\log x}{\log Q}\right)\log \left(\frac{\log x}{\log Q}\right)
%\end{align*}
Since $\frac{5}{8}\beta<\frac{5}{1312}$, we have
$$h(r_2)-g(r_2)\leq -\left(\frac{11}{1312}+o(1)\right)\log \left(u\right).$$
So letting 
$$C_2=\frac{5}{656}\left(\frac{\alpha\kappa \left(1-\frac{\alpha \kappa}{20}\right)}{20}\right),$$
we have
$$\exp\left(-r_2(g(r_2)-h(r_2))\right)\ll \rho(u)^{C_2+o(1)}.$$
Thus,
$$\sum_{n\in IV}f(n)\ll \rho(u)^{C_2+o(1)}\left(\frac{Y}{\phi(k)\log x}\right)\left(\sum_{\substack{J<p\leq Q \\ p\nmid k}}\frac{f(p)}{p}\right).$$
%If $f\in \mathcal M$ as in Shiu's original theorem, we can simply let $J=1$ and not have to worry about the remaining sum over $s$, and thus the above completes the proof.  (In fact, the constant can be improved somewhat, though we do not do this here.)  If $f\in \mathcal(\beta)$, on the other hand, we 
This completes the proof of Theorem \ref{roughbd2}, and by extension, the proof of Theorem \ref{MnThm2}.
%Plugging this into Lemmas SOMETHING and SOMETHING then prove that for any $\epsilon>0$,
%$$\sum_{n\in IV}f(n)\ll \rho(u)^{C_2+o(1)}\left(\frac{Y}{\phi(k)(\log x)^{1-\epsilon}}\right)\left(\sum_{\substack{J<p\leq z^\frac 1r \\ p\nmid k}}\frac{f(p)}{p}\right).$$
%This completes the proof of Theorem \ref{MnThm2}.


%\section{Remark 1: Simplified, Weakened Proof of Theorem \ref{MnThm2}}


%We note that for $f\in \mathcal M_Q$, there is a much simpler way to prove something close to Theorem \ref{MnThm2}, although the simplicity comes at the cost of a slightly worsened bound.  For a given $f\in \mathcal M_Q$, consider the function
%$$j(n)=f(n)n^{\frac{1}{\log Q}}.$$
%Note that if $n\leq Q$ then
%$$j(n)=f(n)e^{\frac{\log n}{\log Q}}\leq ef(n),$$
%while if $n>Q$ then $j(n)=0$.  So Conditions i and ii both apply to $j(n)$, and hence Shiu's theorem applies.  Moreover, by Rankin's trick, we have
%$$\sum_{\substack{x\leq n\leq x+y \\ n\equiv a\pmod k}}f(n)\leq x^{-\frac{1}{\log Q}}\sum_{\substack{x\leq n\leq x+y \\ n\equiv a\pmod k}}j(n).$$
%So by Shiu's theorem,
%$$\sum_{\substack{x\leq n\leq x+y \\ n\equiv a\pmod k}}f(n)\ll \frac{y^{1-\frac{1}{\log Q}}}{\phi(k)(\log x)}\exp\left(e\sum_{\substack{p\leq x \\ p\nmid k}}\frac{f(p)}{p}\right)=\frac{y\rho(u)^{\frac{1+o(1)}{\log u}}}{\phi(k)(\log x)}\exp\left(e\sum_{\substack{p\leq x \\ p\nmid k}}\frac{f(p)}{p}\right).$$

%\section{Remark: Sharpened Theorem \ref{MnThm2} for Characteristic Function on Smooth Numbers}


%We also note that in the case where $f$ is the characteristic function on $Q$-smooth numbers, we can simplify the argument and improve our constant $C_1$ to find the following.  
%\begin{theorem}
%Let $x$, $y$, and $k$ be as in Theorem \ref{Maintheorem}, and let $Z=\frac y{k^{1+\nu}}$ for some positive %$\nu<\frac{\kappa}{2}$.  Let $y=x^\eta$ for some fixed $0<\eta\leq 1$.  Then
%$$\sum_{\substack{x\leq n\leq x+y \\ n\equiv a\pmod k \\ n\in \mathcal S(Q)}}1\ll %\frac{y\rho(u_{x/Z})^{C_1}\rho(u_{Z/Q})\log^2 Q}{\phi(k)\log x}.$$

%\end{theorem}
%\begin{proof}
%Let . For any $n$ that is $Q$-smooth, we can use the same trick as in Lemma \ref{LemmaV3}, writing 
%$$n=cd,$$
%where $\frac ZQ\leq c\leq Z$.  (Again, $c$ and $d$ are not necessarily unique, but this will not concern us.)  We will assume again that $\frac{\log Q}{\log x}=o(1),$ else our bound will be trivial. Note that $$\frac{y+y}{c}\geq \frac %xc+k^{1+\nu}\geq \frac xZ+k^{1+\nu}.$$  So
%\begin{gather}\label{fsmooth}\sum_{\substack{x\leq n\leq x+y \\ n\equiv a\pmod k \\ n\in \mathcal S(Q)}}1\leq \sum_{\substack{\frac ZQ\leq c\leq Z\\ c\in\mathcal S(Q)}}\sum_{\substack{\frac xc\leq d\leq \frac{x+y}{c} \\ d\equiv a\overline{c}\pmod k \\ d\in\mathcal S(Q)}}1\ll \sum_{\substack{\frac ZQ\leq c\leq Z\\ c\in\mathcal S(Q)}}\frac{y\log Q}{c\phi(k)\log x}\rho(u_{x/Z})^{C_1},
%\end{gather}
%where we have applied Theorem \ref{MnThm2} to the sum over $d$.

%We can evaluate the sum over $c$ in much the same way as (\ref{dyad}), finding
%\begin{align*}\sum_{\substack{\frac ZQ\leq c\leq Z\\ c\in\mathcal S(Q)}}&\frac{1}{c}\leq \sum_{j=0}^{\lceil \frac{\log %Q}{\log 2}\rceil}\frac{Q}{2^jZ}\sum_{\substack{2^j\frac ZQ<c\leq 2^{j+1}\frac ZQ \\ c\in \mathcal S(Q)}}1 \\
%\ll &\sum_{j=0}^{\lceil\frac{\log Q}{\log 2}\rceil}\frac{Q}{2^jZ}\left[2^{j+1}\frac ZQ\rho\left(u_{Z/Q}\right)\right]\leq \left(2\log Q\right)\rho\left(u_{Z/Q}\right).
%\end{align*}
%We plug this into (\ref{fsmooth}) to find 
%$$\sum_{\substack{x\leq n\leq x+y \\ n\equiv a\pmod k \\ n\in \mathcal S(Q)}}1\ll %\frac{y\rho(u_{x/Z})^{C_1}\rho(u_{Z/Q})\log^2 Q}{\phi(k)\log x}.$$
%This proves the Theorem.
%\end{proof}
%\begin{corollary}\label{Hilcor}
%Let $Q$, $y$, and $x$ be as in the previous theorem.  Then
%$$\sum_{\substack{x\leq n\leq x+y \\ n\in \mathcal S(Q)}}1\ll y\rho(u_x)^{C_1-(1-C_1)\eta+o(1)}.$$
%Equivalently,
%$$\sum_{\substack{x\leq n\leq x+y \\ n\in \mathcal S(Q)}}1\ll y\rho(u_y)^{1+C_1\left(\frac 1\eta-1\right)+o(1)}.$$
%\end{corollary}

%\begin{proof}
%Let $k=1$ and $Z=y=x^\eta$ in the previous theorem.  We will assume that $\eta$ is a fixed constant less than 1, since %the case of $\eta=1+o(1)$ follows from Hildebrand's result.  We can also again assume that $\frac{\log Q}{\log %x}=o(1),$ else the result is trivial.  So
%\begin{align*}
%\sum_{\substack{x\leq n\leq x+y \\ n\in \mathcal S(Q)}}1\ll &\frac{y\left(\frac{(1-\eta)\log x}{\log %Q}\right)^{-(C_1+o(1))\frac{(1-\eta)\log x}{\log Q}}\left(\frac{\eta\log x}{\log Q}\right)^{-(1+o(1))\frac{\eta\log x}{\log Q}}\log^2 Q}{\log x}\\
%=&y\left(\frac{(1-\eta)\log x}{\log Q}\right)^{-(C_1-(1-C_1)\eta)\frac{\log x}{\log Q}+o(1)}\\
%=&y\rho(u_x)^{C_1-(1-C_1)\eta+o(1)}.
%\end{align*}
%To see that this is less than the Hildebrand bound, we can also write the above as
%$$y\left(\frac{(1-\eta)\log x}{\log Q}\right)^{-(C_1(1-\eta)+\eta)\frac{\log x}{\log Q}+o(1)}\leq y\rho(u_y)^{1+C_1\left(\frac 1\nu-1\right)+o(1)}.$$

%\end{proof}


\section{Application: the $d$-fold Divisor Function}
As in Section 6 of Shiu's original paper, we use this result to prove a statement about the divisor function.  Here, we let $\tau_d$ denote the $d$-fold divisor function, and let $\tau=\tau_2$.
\begin{varthmdfold}
Let $x$, $y$, and $k$ be as in Theorem \ref{Maintheorem}.  Choose an integer $d\geq 2$, and let $R=\frac{\log\log\log x}{g(x)}$ for some $g(x)$ such that $g(x)\to\infty$ as $x\to\infty$.  Then for any $\epsilon>0$,
$$\sum_{\substack{x\leq n\leq x+y \\ n\equiv a\pmod k}}\tau_d(n)^{R}\ll \frac{y(\log x)^\epsilon}{\phi(k)}\left(\frac{\phi(k)}{k}\log x\right)^{d^{R}-1}.$$
Moreover,
$$\sum_{\substack{x\leq n\leq x+y \\ n\equiv a\pmod k \\ n\in \mathcal S(Q)}}\tau(n)^{R}\ll \rho(u)^{C_2+o(1)}\frac{y(\log x)^\epsilon}{\phi(k)}\left(\frac{\phi(k)}{k}\log x\right)^{d^{R}-1}.$$
\end{varthmdfold}
\begin{proof}
For any prime $p$,
$$\tau_d(p)^{R}=d^{\frac{\log\log\log x}{g(x)}}=(\log\log x)^{\frac{\log d}{g(x)}}.$$
So Condition 1 holds for $\beta$ as given above.

Moreover, for any $n\geq 3$, we have an effective upper bound for the binary divisor function \cite{NR}:
$$\tau_2(n)^{R}\leq n^{\frac{1.5379 (\log\log\log x)}{g(x)\log\log n}}$$
Noting that 
$$\tau_{d+1}(n)=\sum_{m|n}\tau_d\left(\frac nm\right)\leq \sum_{m|n}\tau_d(n)=\tau_2(n)\tau_d(n),$$
we can inductively bound $\tau_d$ by
$$\tau_d(n)\leq \tau_2(n)^{d-1}.$$
Hence
$$\tau_d(n)^{R}\leq n^{\frac{1.5379 (d-1)(\log\log\log x)}{g(x)\log\log n}}$$
The right-hand side is obviously an increasing function on $n$.  So if $n\leq (\log x)^{10}$ then this is 
$$\leq (\log x)^{\frac{15.379 (d-1)}{g(x)}},$$
which is clearly $\ll (\log x)^\epsilon$ for any $\epsilon$.  If $n>(\log x)^{10}$ then this is
$$\leq n^{\frac{1.5379  (d-1)}{g(x)}},$$
which is $\ll n^\epsilon$ for every $\epsilon>0$.  So Condition 2 also holds.  Hence, $\tau_d(n)^R\in M(\beta)$, where $\beta$ is as in Theorem \ref{Maintheorem}.

Applying Theorem \ref{Maintheorem}, we then have
$$\sum_{\substack{x\leq n\leq x+y \\ n\equiv a\pmod k}}\tau_d(n)^{R}\ll \frac{y}{\phi(k)(\log x)^{1-\epsilon}}\exp\left(d^R\sum_{\substack{p\leq x \\ p\nmid k}}\frac{1}{p}\right)\ll \frac{y(\log x)^\epsilon}{k} \left(\frac{\phi(k)}{k}\log x\right)^{d^{R}-1}$$
as in the original proof in \cite{Sh}.

The proof of the second half of the theorem is nearly identical except that we apply Theorem \ref{MnThm2} instead of Theorem \ref{Maintheorem}.

\end{proof}

%By similar process, we can find another result over smooth numbers:
%\begin{theorem}\label{divsmooth}
%Let $x$, $y$, $k$, and $\beta$ be as in Theorem \ref{Maintheorem}.  Choose an integer $d\geq 2$, and let %$R=\frac{\log\log\log x}{g(x)}$ for some $g(x)$ such that $g(x)\to\infty$ as $x\to\infty$.  Then for any $\epsilon>0$,
%$$\sum_{\substack{x\leq n\leq x+y \\ n\equiv a\pmod k \\ n\in \mathcal S(Q)}}\tau(n)^{R}\ll \rho(u)^{C_2+o(1)}\frac{y(\log x)^\epsilon}{\phi(k)}\left(\frac{\phi(k)}{k}\log x\right)^{d^{R}-1}.$$
%\end{theorem}
`%\begin{proof}
%Applying Theorem \ref{MnThm2} instead of Theorem \ref{Maintheorem} gives the above result.
%\begin{align*}
%\sum_{\substack{x\leq n\leq x+y \\ n\equiv a\pmod k}}\tau_d(n)^{R}\ll &\frac{y}{\phi(k)(\log %x)^{1-\epsilon}}\exp\left(d^R\sum_{\substack{p\leq Q \\ p\nmid k}}\frac{1}{p}\right)\\
%=&\frac{y}{\phi(k)(\log x)^{1-\epsilon}}\exp\left(d^R\sum_{\substack{p\leq Q }}\frac{1}{p}-d^R\sum_{\substack{p\leq Q %\\ p|k}}\frac 1p\right)\\
%=&\frac{y}{\phi(k)(\log x)^{1-\epsilon}}\left(\frac{\phi(k)\log Q}{k}\right)^{d^R}\exp\left(d^R\sum_{\substack{p>Q \\ p|k}}\frac 1p\right)\\
%\ll \frac{x(\log x)^\epsilon}{k} \left(\frac{\phi(k)}{k}\log x\right)^{d^{R}-1}
%\end{align*}
%\end{proof}
We note that if we were to prove the smooth portion of Theorem \ref{dfold} by restricting our sum over the primes in the exponential term to $\sum_{p\leq Q}\frac{\tau_d(p)}p$, we would only save a term that is $\rho(u)^{o(1)}$, and hence we do not pursue this idea further here.

\section{Application: Smooth Numbers in Short Intervals}
Finally, we turn to the results of \cite{So} and \cite{Go}.  We find the following improvement.


\begin{vtsmooththm}
Let $Q$, $x$, and $y$ satisfy either the hypotheses in either Theorem \ref{SoundThm} or \ref{GoThm}, where in the former case we also assume the Riemann Hypothesis.  Additionally, assume that $u\geq (\log\log x)^2.$  Then there exists a constant $C_3$ such that 
\begin{gather}
\psi(x+y,Q)-\psi(x,Q)\gg y\rho(u)^{1+\frac{C_3\log\log \log\log x}{\log\log \log x}}.
\end{gather}
In particular, we can take $C_3>1-C_2$.
\end{vtsmooththm}


%\begin{theorem}
%Assume the Riemann Hypothesis, and choose $y$ such that $u\geq \left(\log\log x\right)^2.$  Then there exists a $C_3$ %such that 
%$$y\rho(u)^{1-\frac{C_3\log\log\log \log x}{\log\log\log x}}\ll \sum_{\substack{x\leq n\leq x+y \\ n\in \mathcal S(Q) %}}1.$$

%\end{theorem}
\begin{proof}
In both of these cases, the first step is to find a lower bound for $$\sum_{\substack{x\leq n\leq x+y \\ n\in \mathcal S(Q) }}\tau_d(n)$$
for some $d$ and for $y$ as in the theorem.

We begin with \cite{So}.  Using the notation of that paper, we define $\delta$ such that
$$xe^{2\delta}=x+y.$$
Naturally, this means that $$\delta=\frac{y}{2x}(1+o(1)).$$
The paper also defines
$$I=\sum_{x\leq n\leq x+y}\sum_{\substack{n=rm_1m_2 \\ m_1,m_2\in \mathcal S(Q) \\ \sqrt xQ^{-\frac 13}\leq m_1,m_2\leq xQ^{-\frac 14} }}\Lambda(r)\min\{\log \frac{e^{2\delta}x}{n},\log \frac nx\}.$$
Note that $r\leq Q$, and hence $r$ is also $Q$-smooth.  So
$$I\ll \delta\log Q\sum_{\substack{x\leq n\leq x+y \\ n\in \mathcal S(Q) }}\tau_3(n).$$
This paper eventually finds that 
$$I\gg x\delta^2M(1)^2/2,$$
where $M(1)\geq \rho(u/2)\log Q/24$.
Putting these together, we have
$$x\delta\rho(u/2)^2\log Q\ll \log Q\sum_{\substack{x\leq n\leq x+y \\ n\in \mathcal S(Q) }}\tau_3(n),$$
which can be re-expressed as
\begin{gather}\label{zrho}y\rho(u/2)^2\ll \sum_{\substack{x\leq n\leq x+y \\ n\in \mathcal S(Q) }}\tau_3(n).
\end{gather}
Moving on to the unconditional case of \cite{Go}, the process is similar.  Near the end of Section 6 of \cite{Go}, the author of that paper finds that 
$$\rho(u/2)^2h\sqrt x\ll \sum_{\substack{x-2h\sqrt x<n_1n_2\leq x+2h\sqrt x \\ n_1\sim \sqrt x \\ n_1,n_2\in \mathcal S(Q)}}1.$$
Clearly,
$$\sum_{\substack{x-2h\sqrt x<n_1n_2\leq x+2h\sqrt x \\ n_1\sim \sqrt x \\ n_1,n_2\in \mathcal S(Q)}}1\ll \sum_{\substack{x-2h\sqrt x<n\leq x+2h\sqrt x \\ n\in \mathcal S(Q)}}\tau(n),$$
and hence if we rewrite $x-2h\sqrt x$ as $x$ and let $y=4h\sqrt x$ then we have
\begin{gather}\label{zrho2}y\rho(u/2)^2\ll \sum_{\substack{x<n\leq x+y \\ n\in \mathcal S(Q)}}\tau(n).\end{gather}
Now, let 
$$R=\frac{\log\log \log x}{\log \log \log \log x}.$$
Then by Holder's inequality,
$$\sum_{\substack{x\leq n\leq x+y \\ n\in \mathcal S(Q) }}\tau_d(n)\leq \left(\sum_{\substack{x\leq n\leq x+y \\ n\in \mathcal S(Q) }}1^{\frac{1}{1-\frac 1R}}\right)^{1-\frac 1R}\left(\sum_{\substack{x\leq n\leq x+y \\ n\in \mathcal S(Q) }}\tau_d(n)^{R}\right)^\frac 1R.$$
Applying Theorem \ref{dfold}, we that the above is
\begin{align*}\ll &\left(\sum_{\substack{x\leq n\leq x+y \\ n\in \mathcal S(Q) }}1\right)^{1-\frac 1R}\left(y \rho(u)^{C_2+o(1)}\left(\log x\right)^{d^{R}-1+\epsilon}\right)^\frac 1R\\
\ll &\left(\sum_{\substack{x\leq n\leq x+y \\ n\in \mathcal S(Q) }}1\right)^{1-\frac 1R}y^\frac 1R \rho(u)^{\frac{\log\log\log\log x}{\log\log\log x}\left(C_2+o(1)\right)}\left(\log x\right)^{\left(\log\log x\right)^{\frac{\log d}{\log\log\log\log x}}}\\
\ll & \left(\sum_{\substack{x\leq n\leq x+y \\ n\in \mathcal S(Q) }}1\right)^{1-\frac 1R}y^\frac 1R \rho(u)^{\frac{\log\log\log\log x}{\log\log\log x}\left(C_2+o(1)\right)}e^{u^\frac 23}.
\end{align*}
since $u\geq (\log\log x)^2$. 

Recall that
$$\rho(u)=\exp\left(-u\left(\log u+\log\log u-1+\frac{\log\log u-1}{\log u}+O\left(\left(\frac{\log\log u}{\log u}\right)^2\right)\right)\right).$$
So
%$$\rho(u/2)=\exp\left(-\frac u2\left(\log u-\log 2+\log\log u-\frac{\log 2}{\log u}-1+\frac{\log\log u-\frac{\log 2}{\log u}-1}{\log u}+O\left(\left(\frac{\log\log u}{\log u}\right)^2\right)\right)\right).$$
%Hence,
$$\rho(u/2)^2=\exp\left(-u\left(\log u-\log 2+\log\log u-1+O\left(\frac{\log\log u}{\log u}\right)\right)\right).$$
Plugging into (\ref{zrho}) or (\ref{zrho2}), we then have
\begin{align*}
y\exp\left(-u\left(\log u+\log\log u+O(1)\right)\right)\ll \left(\sum_{\substack{x\leq n\leq x+y \\ n\in \mathcal S(Q) }}1\right)^{1-\frac 1R}y^\frac 1R\rho(u)^{\frac{\log\log\log\log x}{\log\log\log x}\left(C_2+o(1)\right)}e^{u^\frac 23},
\end{align*}
or
\begin{align*}
y\left[\exp\left(-\left(1-\frac{C_2}{R}\right)u\left(\log u+\log\log u+O(1)\right)\right)\right]^{1+\frac{1}{R}}\ll \sum_{\substack{x\leq n\leq x+y \\ n\in \mathcal S(Q) }}1,
\end{align*}
Distributing the exponent gives us a left-hand side of
$$y\exp\left(-\left(1+\frac{1-C_2}{R}-\frac{C_2}{R^2}\right)u\left(\log u+\log\log u+O(1)\right)\right).$$
Plugging back in for $R$ gives
\begin{align*}
y\rho(u)^{1+\frac{(1-C_2)\log\log\log \log x}{\log\log\log x}+O\left(\left(\frac{\log\log \log u}{\log \log \log \log u}\right)^2\right)}\ll \sum_{\substack{x\leq n\leq x+y \\ n\in \mathcal S(Q) }}1,
\end{align*}
which is as required.
\end{proof}
\section{Acknowledgements}

I would like to thank Sarvagya Jain for helpful comments and suggestions.

\bibliographystyle{line}
%\bibliography{JAMS-paper}
\begin{thebibliography}{[HD]}
\normalsize
\bibitem{Di} K. Dickman, On the frequency of numbers containing prime factors
of a certain relative magnitude, Ark. Mat. Astr. Fys. 22 (1930), 1-14.
\bibitem{Gr} A. Granville, Smooth numbers: computational number theory and beyond, MSRI Publ., 44 (2008), 267-323.
\bibitem{Go}  \'{E}. Goudout, Lois locales de la fonction $\omega$ dans presque tous les petits intervalles, Proc. Lond. Math. Soc. (3) 115(3) (2017), 599-–637.
\bibitem{Hi85} A. Hildebrand, Integers free of large prime divisors in short
intervals, Quart. J. Math. Oxford Ser. 2, 36:141 (1985), 57–69.
\bibitem{Hi}  A. Hildebrand, On the number of positive integers $\leq x$ and free of prime factors $>y$, J. Number Theory 22:3 (1986), 289-307.
%\bibitem[Ma]{Ma}  K. Matom\"{a}ki. Another note on smooth numbers in short intervals. Int. J. Number Theory, 12(2) (2016) 323–-340.
\bibitem{Ja} S. Jain, Smooth Numbers in Short Intervals, https://arxiv.org/pdf/2502.10530.
\bibitem{NR} J.L. Nicolas and G. Robin, Majorations Explicites Pour le Nombre de Diviseurs de N, Can. Math. Bull. 26(4) (1983), 485--492.
\bibitem{Sh} P. Shiu, A Brun-Titchmarsh theorem for multiplicative functions, J. Reine Angew. Math. 313
(1980), 161-170.
\bibitem{So}  K. Soundararajan. Smooth numbers in short intervals, https://arxiv.org/pdf/1009.1591v1.
\end{thebibliography}


\end{document} 
